\section{Removing rough progression steps}
\label{sec:rough-progressions}

At an alignment pivot $i$, the input of a model has the form $t_Tt_i$, where
the earlier variables in $t_T$ are temporarily fixed.  Consequently,
averaging $t_i$ samples only one progression of step $t_T$ in the
model input.  We will need to remove the factors of this step while
retaining polynomial dependence on the factor being removed.  The
comparison proved here permits the sampling box, residue and observable
to depend arbitrarily on that factor.  This last uniformity is necessary
because the actual pivot cell is determined only after sampling.

All nilmanifolds in this Section are quotients $G/\Gamma$, with $G$
connected, simply connected and nilpotent and $\Gamma$ a lattice.  We fix
rational coordinates on $\mathfrak g=\log G$ and a smooth metric on the
compact quotient.  ``Polynomial in logarithmic coordinates'' means that
the coordinates of $\log P$ are ordinary real polynomials.  Their degrees,
but not their coefficients, will be bounded.  For a scalar $a$, write
$\norm{a}_{\R/\Z}=\inf_{k\in\Z}\abs{a-k}$.

The orbit theory underlying this comparison begins with Leibman's
qualitative equidistribution theorem \cite{LeibmanOrbits}.
We use the quantitative form of Green and Tao, with its multiparameter
correction, in the precise version stated next.

\subsection{The corrected quantitative equidistribution input}

A rational filtration is a finite decreasing sequence of connected
rational subgroups $G_\bullet=(G_j)_{j\geq0}$ with
$G_0=G_1=G$ and $[G_i,G_j]\subseteq G_{i+j}$.  A map on $\Z^d$ is
polynomial for this filtration if every $k$-fold group difference takes
values in $G_k$.  Fix a rational Mal'cev basis adapted to the filtration.
A horizontal character is a continuous homomorphism
$\xi:G\longrightarrow\R$ such that $\xi(\Gamma)\subseteq\Z$; its
size is the norm of its integer coefficient vector on the horizontal
torus.  In particular, characters of bounded size form a finite set.

\begin{theorem}[Quantitative Leibman theorem, corrected box form]
\label{thm:quantitative-leibman}
Fix a rational filtered nilmanifold as above, a number $d$ of variables,
and $0<\delta<1$.  There is a constant $A$, depending only on these data,
with the following property.  If $P:\Z^d\to G$ is polynomial for this
filtration and $(P(v)\Gamma)_{v\in\{0,\ldots,N-1\}^d}$ is not
$\delta$-equidistributed, then there is a nonzero horizontal character
$\xi$ of size at most $A$ such that, on writing
\[
 \xi(P(v))=\sum_I\alpha_I\binom vI,
 \qquad \binom vI=\prod_{j=1}^d\binom{v_j}{I_j},
\]
one has
\begin{equation}
 N^{\abs I}\norm{\alpha_I}_{\R/\Z}\leq A
 \qquad (\abs I>0).
 \label{eq:rough-leibman-binomial}
\end{equation}
Here $\delta$-equidistribution means comparison with Haar probability to
error at most $\delta$ against every observable of Lipschitz norm at most
one.  The constant $A$ is independent of the coefficients of $P$ and of
$N$.

The following consequence will be used.  Fix also $c,C,B,\eta>0$.  Let
$\mathcal Q\subseteq[-CZ,CZ]^d$ be an axis-parallel box with each side
at least $cZ$, sampled at its integer points.  If
\begin{equation}
 \abs{\E_{x\in\mathcal Q\cap\Z^d}F(P(x)\Gamma)
       -\int_{G/\Gamma}F}\geq\eta,
 \qquad \norm F_{\Lip}\leq B,
 \label{eq:rough-haar-obstruction}
\end{equation}
then, for sufficiently large $Z$, there is a nonzero horizontal
character in a fixed finite list such that, writing
$\xi(P(x))=\sum_I\theta_Ix^I$ in ordinary monomials,
\begin{equation}
 \norm{\theta_I}_{\R/\Z}\leq A'Z^{-\abs I}
 \qquad (\abs I>0).
 \label{eq:rough-leibman-monomial}
\end{equation}
The list and $A'$ depend only on the fixed filtered nilmanifold,
$d,c,C,B,\eta$.  The same assertions hold with the nilmanifold and its
filtration chosen from a fixed finite list.
\end{theorem}

The equal-side statement is the corrected version of
\cite[Theorem~8.6]{GreenTao}; see the expanded erratum
\cite[p.~3 and Section~3, arXiv:1311.6170v3]{GreenTaoErratum}.
The one-variable case is
\cite[Theorem~2.9]{GreenTao}, which is unaffected by the correction.
The published quantitative statement bounds the rationality of the
adapted basis by the reciprocal accuracy.  Here the basis is fixed;
decreasing the accuracy below the reciprocal of its fixed rationality
bound absorbs that hypothesis into $A$.
We use the equal-side case of Theorem~\ref{thm:quantitative-leibman}
and now verify the stated box consequence;
no assertion about unrestricted unequal side lengths is needed.

\begin{proof}[Derivation of the box consequence]
Normalize the observable by its Lipschitz bound.  Choose a small positive
$\lambda$, depending only on $d,c,C,B,\eta$, and tile the integer box by
cubes of side $\floor{\lambda Z}$, leaving strips along its boundary.
The strips contain at most an $O_{d,c}(\lambda+Z^{-1})$ fraction of its
integer points.  Choose $\lambda$ so that their contribution to
\eqref{eq:rough-haar-obstruction} is less than $\eta/4$.  Expressing
the remaining average as a convex combination shows that one cube has
Haar discrepancy at least $\eta/2$. Apply the equal-side case of
Theorem~\ref{thm:quantitative-leibman} to
$P(a+v)$ on this cube.  Integer translations preserve filtered
polynomiality because a group difference of a translated map is the
corresponding translate of its group difference.

There is a fixed bound $D_0$ on the scalar degree.  Multiply the resulting
character by $(D_0!)^d$.  Indeed, on expanding the products
$\binom vI$, this multiplication clears every denominator, so each
ordinary monomial coefficient is an integer plus an error
$O(Z^{-\abs I})$: a coefficient of degree $I$ receives contributions
only from binomial degrees $J\geq I$, and
$Z^{-\abs J}\leq Z^{-\abs I}$.  Finally expand $(x-a)^J$.
The translation has integer entries and $\abs a=O_{C,d}(Z)$; hence
its integer coefficient contributions stay integral, while its errors
at degree $I$ are bounded by
\[
 \sum_{J\geq I}O(Z^{\abs J-\abs I})O(Z^{-\abs J})
 =O(Z^{-\abs I}).
\]
This proves \eqref{eq:rough-leibman-monomial}.  The bounded multiplication
of the character still leaves a finite list.  The bound on the location
of $\mathcal Q$ is used precisely in this last translation estimate.
\end{proof}

We explain why Theorem~\ref{thm:quantitative-leibman} applies to the logarithmic polynomials used
below.  Let $G^{(j)}$ be the lower central series, with $G^{(1)}=G$, and
suppose $G$ has step at most $s$.  For an ordinary degree bound $D\geq1$,
put
\begin{equation}
 G_0=G,\qquad G_i=G^{(\ceil{i/D})}\quad(i\geq1).
 \label{eq:rough-stretched-filtration}
\end{equation}
These are connected rational subgroups, the filtration ends after $sD$,
and
$\ceil{i/D}+\ceil{j/D}\geq\ceil{(i+j)/D}$ gives the bracket
condition; the condition with $i=0$ follows because the lower central
subgroups are normal.  If $\log P$ has ordinary total degree at most
$D$, its coefficient at a monomial of degree $j$ belongs to
$\mathfrak g_j$, since $\mathfrak g_j=\mathfrak g$ for $j\leq D$.

Here is a direct verification of polynomiality, including the effect
of noncommutativity.  Suppose the coefficient at degree $j$ in
$A(x)=\log Q(x)$ lies in $\mathfrak g_{j+k}$.  An ordinary difference
$A(x+h)-A(x)$ has its degree-$j$ coefficient in
$\mathfrak g_{j+k+1}$, because it arises from degrees at least $j+1$.
Using the Baker--Campbell--Hausdorff (BCH) formula, expand
\[
 \log\bigl(Q(x+h)Q(x)^{-1}\bigr)
   =\operatorname{BCH}\bigl(A(x)+[A(x+h)-A(x)],-A(x)\bigr).
\]
The terms containing no occurrence of $A(x+h)-A(x)$ cancel, since
$\operatorname{BCH}(A,-A)=0$.  Every remaining bracket contains at
least one difference.  Filtration indices add in brackets, so its
coefficient at degree $j$ belongs to $\mathfrak g_{j+k+1}$.
The BCH series is finite.  Induction starting at $k=0$ proves that the
$k$-fold group difference takes values in $G_k$.  Its degree is bounded
in terms of $s,D$ throughout.  Thus arbitrary coefficients in a
bounded-degree logarithmic polynomial are allowed in
Theorem~\ref{thm:quantitative-leibman}, using the fixed filtration
\eqref{eq:rough-stretched-filtration}.  The same argument works with
any fixed number of ordinary variables.

\subsection{An interpolation fact with smooth denominators}

The arithmetic reason that a rough step can be removed is the following
fact.  Its analytic constants do not depend on the length of the
progression.  Rational denominators may depend on that length, which is
always fixed before taking the asymptotic limit.

\begin{lemma}[Exact division after interpolation]
\label{lem:rough-interpolation}
Fix $e\geq1$, $q\geq e$, and $A>0$.  For a fixed integer
$H\geq2q+2$, let
\[
 t_z=t_0+mz\in[S,2S),\qquad 0\leq z<H,
\]
where $m$ is a positive integer with all prime factors at most $w$, and
$\gcd(t_z,W)=1$ for every $z$, with $W=\prod_{p\leq w}p$.
Suppose $w,S\to\infty$ and $Z/S^a\to\infty$ for every fixed $a>0$.
Let $\theta(z)$ be a real polynomial such that
$\deg(t_z^e\theta(z))\leq q$ and
\begin{equation}
 \norm{t_z^e\theta(z)}_{\R/\Z}\leq A(S/Z)^e
 \qquad(0\leq z<H).
 \label{eq:rough-interpolation-input}
\end{equation}
For all sufficiently large $w$ there is a rational polynomial $M(z)$,
of degree at most $q-e$, integral at every index $0\leq z<H$, such that
\begin{equation}
 \abs{\theta(z)-M(z)}\leq C_{q}A Z^{-e}
 \qquad(0\leq z<H).
 \label{eq:rough-interpolation-output}
\end{equation}
Its coefficient denominators have only prime factors at most $w$.
The constant in \eqref{eq:rough-interpolation-output} is independent of
$H$.  The threshold for $w$ is permitted to depend on $H$.
\end{lemma}

\begin{proof}
Write $T(z)=(t_0+mz)^e\theta(z)$ and
$\epsilon=A(S/Z)^e$.  Choose $q+1$ integer nodes
$z_0,\ldots,z_q$ in $[0,H-1]$ with pairwise successive gaps at least
$(H-1)/(2q)$, for example the nearest integers to $j(H-1)/q$.
Let $n_j$ be a nearest integer to $T(z_j)$, and let $R$ be their
degree-at-most-$q$ Lagrange interpolant.  On rescaling the index interval
to $[0,1]$, the interpolation nodes remain separated by $1/(2q)$.
The Lagrange basis polynomials and all their coefficients in the
rescaled variable are therefore bounded in terms of $q$ alone.  Since
$T$ already has degree at most $q$, interpolation of the errors gives
\begin{equation}
 \sup_{0\leq z\leq H-1}\abs{R(z)-T(z)}\leq C_q\epsilon.
 \label{eq:rough-interpolation-uniform}
\end{equation}
There is a positive integer $D_H$, depending only on the chosen nodes,
such that $D_HR\in\Z[z]$.  One may take the product of the nonzero
node differences occurring in the Lagrange denominators.  In particular,
$D_H$ is independent of $w,t_0,m$ and the polynomial coefficients.

Let $a=-t_0/m$.  The polynomial $T$ has a zero of multiplicity at least
$e$ at $a$.  The coefficient bounds in the rescaled Lagrange formula
give, for $0\leq j<e$,
\[
 \abs{R^{(j)}(a)}
  =\abs{(R-T)^{(j)}(a)}
  \leq C_q\epsilon(H-1)^{-j}
          \left(1+\frac{\abs a}{H-1}\right)^q
  \leq C_q\epsilon(1+2S)^q.
\]
On the other hand $R^{(j)}(a)$ is rational with denominator dividing
$D_Hm^q$.  We have $m\leq S$ because the progression lies in $[S,2S)$
and has at least two terms.  Consequently a nonzero such rational has
absolute value at least $D_H^{-1}S^{-q}$.  For fixed $H$, the displayed
upper bound is smaller than this number for large $w$, by the assumed
domination of every power of $S$ by $Z$.  Thus
$R^{(j)}(a)=0$ for all $j<e$, exactly.  Polynomial division now gives
\begin{equation}
 R(z)=(t_0+mz)^e M(z),\qquad M\in\Q[z],\qquad\deg M\leq q-e.
 \label{eq:rough-exact-divisibility}
\end{equation}
In division by $(t_0+mz)^e$, the only denominators introduced, besides
those already dividing $D_H$, come from its leading coefficient $m^e$.
All their prime factors are therefore at most $w$ once $w$ exceeds the
prime factors of the fixed integer $D_H$.

For every integer $z\in[0,H-1]$, \eqref{eq:rough-interpolation-input}
and \eqref{eq:rough-interpolation-uniform} place $R(z)$ within
$(C_q+1)\epsilon$ of an integer.  Since $R(z)\in D_H^{-1}\Z$, it is
itself an integer for large $w$.  Write $M(z)=a_z/b_z$ in lowest terms.
Every prime factor of $b_z$ is at most $w$, whereas
$\gcd(t_z,W)=1$.  The integrality of $t_z^eM(z)=R(z)$ then forces
$b_z=1$.  Finally divide \eqref{eq:rough-interpolation-uniform} by
$t_z^e\geq S^e$.  This yields
\eqref{eq:rough-interpolation-output} with a constant depending on $q$
alone.  Only the rational exactness steps, not this constant, required
a threshold depending on $H$.
\end{proof}

\subsection{Uniform progression comparison}

Call an integer $w$-smooth if all its prime factors are at most $w$.
For a nonempty finite set $A$ we use $\E_A$ for normalized counting.
The next Proposition compares two sampling laws for the same polynomial
family.  Its joint polynomial dependence on $t$ and $x$ is essential;
the box, residue class, and observable may be chosen separately for
each $t$.

\begin{proposition}[Removal of a rough step]
\label{prop:rough-step}
Fix $G/\Gamma$, a number $d\geq1$ of spatial variables, a degree bound
$D$, and constants $c,C,B>0$.  Along an arbitrary sequence $w\to\infty$,
let $L,S,Z$ satisfy
\begin{equation}
 W\mid L,\quad L\text{ is }w\text{-smooth},\quad
 S\to\infty,\quad S/L\to\infty,\quad
 Z/S^a\to\infty\quad\text{for every fixed }a>0.
 \label{eq:rough-scales}
\end{equation}
Fix a residue $r$ modulo $L$ with $\gcd(r,W)=1$, and put
$\mathcal T=[S,2S)\cap(r+L\Z)$.
Let $P_t(x)$ be $G$-valued, with logarithmic coordinates polynomial
jointly in $(t,x)\in\R\times\R^d$ of total degree at most $D$.
These polynomials and their coefficients may change with $w$.

For every $\eta>0$, the proportion of $t\in\mathcal T$ for which
there exist a box $\mathcal Q\subseteq[-CZ,CZ]^d$ with all sides at
least $cZ$, a vector $b\in\Z^d$, and an observable
$F:G/\Gamma\to\C$ with $\norm F_{\Lip}\leq B$, such that
\begin{equation}
 \abs{\E_{x\in\mathcal Q\cap\Z^d}F(P_t(x)\Gamma)
       -\E_{x\in\mathcal Q\cap(b+t\Z^d)}F(P_t(x)\Gamma)}
       \geq\eta
 \label{eq:rough-discrepancy}
\end{equation}
tends to zero.

This conclusion holds for every sequence of the allowed data.
In particular the exceptional proportion is uniform over the
polynomial coefficients, the residue $r$, and all the boxes, residue
vectors and observables in \eqref{eq:rough-discrepancy}.  The latter
three may be chosen separately for every $t$ without any polynomial
dependence.  The same conclusion holds for a fixed finite list of
nilmanifolds and fixed complexity bounds.
\end{proposition}

\begin{proof}
The induction will turn a discrepancy on $G/\Gamma$ into one on the
kernel of a horizontal character.  To preserve polynomial dependence
during this reduction, we work on long progressions of exceptional
values of $t$ with smooth difference.  Interpolation on such a
progression will split the horizontal polynomial into an integer part
and a slowly varying error.  We first specify the stronger progression
assertion to which this dimension induction applies.

A \emph{bad progression family} consists of fixed
bounds $d,D,c,C,B,\eta>0$ and a fixed nilmanifold such that, for every
prescribed integer $H$, there are sequences with $w,S\to\infty$ and
$Z/S^a\to\infty$ for every fixed $a$, and progressions
\begin{equation}
 t_z=t_0+mz\in[S,2S),\quad 0\leq z<H,\quad
 m\text{ is }w\text{-smooth},\quad \gcd(t_z,W)=1,
 \label{eq:rough-bad-progression}
\end{equation}
on which every index has a discrepancy at least $\eta$ as in
\eqref{eq:rough-discrepancy}.  Here the logarithms of $P_z(x)$ need
only be jointly polynomial in $(z,x)$ with the fixed degree bound;
the boxes, classes and tests may vary with $z$.  The sequences may
depend on $H$.  It is equivalent to ask for arbitrarily large prescribed
lengths: a longer progression can be restricted.  We will prove that
no such family exists, by induction on $\dim G$.

First, failure of the Proposition would produce such a family.  Indeed,
suppose the proportion of bad values is at least $\alpha>0$ along a
subsequence.  In the index coordinate $t=r+Lu$, the set $\mathcal T$
is an integer interval of length tending to infinity.  For prescribed
$H$, finite Szemer\'edi's Theorem \cite{Szemeredi} supplies a fixed
integer $K=K(H,\alpha)$ such that every subset of $[K]$ of density
at least $\alpha/2$ contains an $H$-term progression.  Partition the
index interval into blocks of length $K$ and discard its final
incomplete block.  Some full block has bad density at least
$\alpha/2$ for large $w$.  Its bad progression has step $jL$ with
$1\leq j\leq K$.  This step is $w$-smooth once $w\geq K$.
Joint polynomiality is preserved by the affine substitution for $t$.
Thus \eqref{eq:rough-bad-progression} holds for each fixed $H$.

We shall repeatedly thin a bad progression to one color in a finite
coloring.  Finite van der Waerden's Theorem \cite{vanderWaerden} makes
this legitimate: to obtain a length $H$ of one color, start with the
fixed length prescribed by that Theorem for $H$ and the number of
colors.  Under $z=a+bu$, the new step is $bm$; $b$ is bounded in terms
of the initial fixed length, so it remains $w$-smooth for large $w$.
The new step is still at most $S$, and all scale and degree hypotheses
remain valid.  A finite choice of group data can also be fixed: first
pass to a subsequence for each length, then choose one of the finitely
many possibilities occurring for unbounded lengths.  Bounds in all
these colorings will be independent of $H$.

For dimension zero the quotient is a point and every discrepancy is
zero.  Assume the assertion has been proved in smaller dimension, and
suppose that a bad progression family on $G/\Gamma$ exists.
For each $z$, at least one of its two samples has Haar discrepancy at
least $\eta/2$.  Write its step as $p_z\in\{1,t_z\}$ and its
residue representative as $b_z\in\{0,\ldots,p_z-1\}^d$.
Substitute $x=p_zv+b_z$.  The resulting box lies in a fixed multiple
of $[-Z/p_z,Z/p_z]^d$ and has comparable sides.  Its common scale
tends to infinity.  Theorem~\ref{thm:quantitative-leibman}, with the
fixed stretched filtration \eqref{eq:rough-stretched-filtration},
therefore gives a character from a fixed finite list.  Thin to a common
nonzero character $\xi$, including the fixed denominator multiplication
in that choice.  Write
\begin{equation}
 \xi(P_z(x))=\sum_I\theta_I(z)x^I.
 \label{eq:rough-character-polynomial}
\end{equation}
The coefficients $\theta_I$ are polynomials of bounded degree in $z$.
The coefficient obstruction in the $v$ variables is
\begin{equation}
 \left\|p_z^{\abs I}
       \sum_{J\geq I}\binom JI b_z^{J-I}\theta_J(z)
       \right\|_{\R/\Z}
       \leq A_0(p_z/Z)^{\abs I}
 \qquad(\abs I>0),
 \label{eq:rough-substituted-obstruction}
\end{equation}
where $A_0$ is independent of $H,z$ and of every polynomial coefficient.

We claim, descending through the positive spatial degrees, that there
are rational polynomials $m_I(z)$ of bounded degree such that
\begin{equation}
 m_I(z)\in\Z,\qquad
 \abs{\theta_I(z)-m_I(z)}\leq A_1 Z^{-\abs I}
 \quad(0\leq z<H,\ \abs I>0),
 \label{eq:rough-integer-coefficients}
\end{equation}
with smooth coefficient denominators and a constant $A_1$ independent
of $H$.  Take $H$ larger than the finitely many degree requirements of
Lemma~\ref{lem:rough-interpolation}; this does not restrict the
arbitrarily long progression assertion.  Suppose higher degrees have
already been treated and put $e=\abs I$.  Their integral parts in
\eqref{eq:rough-substituted-obstruction} contribute integers.  Their
error contribution has size at most
\[
 C\sum_{J>I}p_z^e p_z^{\abs J-e}Z^{-\abs J}
 \leq C'(p_z/Z)^e,
\]
since $p_z\leq2S=o(Z)$; here $J>I$ means $J\geq I$ and $J\ne I$.
It follows that
$\norm{p_z^e\theta_I(z)}_{\R/\Z}\leq C'(p_z/Z)^e$.
If $p_z=t_z$ this is already the needed bound.  If $p_z=1$,
multiplication by the integer $t_z^e$ gives the same conclusion:
\[
 \norm{t_z^e\theta_I(z)}_{\R/\Z}\leq C''(S/Z)^e.
\]
Lemma~\ref{lem:rough-interpolation} gives $m_I$ and the claimed bound.
There are only boundedly many spatial degrees and coefficients, so
the descending induction terminates with $A_1$ depending only on the
fixed data.  The length-dependent thresholds for rational exactness
can be met simultaneously.

The positive-degree coefficients now have the required integer parts
and small errors.  To obtain a bounded horizontal error on the whole
box, we must also split the spatially constant coefficient.  This
requires a finite coloring argument.  Let $q_0$ bound
$\deg\theta_0$.  Color $z$ by
the interval containing $\{\theta_0(z)\}$ in a partition of $[0,1)$
into intervals of length less than $2^{-q_0-2}$.  Thin to a long
monochromatic progression and reparametrize it.  On each consecutive
$(q_0+2)$-term segment of the new progression,
\[
 \Delta^{q_0+1}\floor{\theta_0(z)}
       =-\Delta^{q_0+1}\{\theta_0(z)\}
\]
is an integer of absolute value less than one: the constant part of
the common fractional-part interval cancels, and the sum of the
absolute difference coefficients is $2^{q_0+1}$.  It is therefore zero.
Newton interpolation, or induction using this difference identity,
shows that a rational polynomial $m_0(z)$ of degree at most $q_0$
agrees with $\floor{\theta_0(z)}$ on the whole new progression.
Its Newton coefficients are integer differences; hence its ordinary
coefficient denominators divide $q_0!$.  All previously constructed
$m_I$ retain \eqref{eq:rough-integer-coefficients} after the same
affine reparametrization.

Set
\begin{equation}
 m_z(x)=\sum_I m_I(z)x^I,
 \qquad a_z(x)=\xi(P_z(x))-m_z(x).
 \label{eq:rough-horizontal-splitting}
\end{equation}
For integer $z$ in the progression, $m_z$ has integer coefficients
as a polynomial in $x$.  By \eqref{eq:rough-integer-coefficients},
its error satisfies, throughout $[-CZ,CZ]^d$,
\begin{equation}
 \abs{a_z(x)}\leq C_1,
 \abs{\partial_{x_j}a_z(x)}\leq C_1/Z
 \quad(1\leq j\leq d),
 \label{eq:rough-slow-factor}
\end{equation}
where $C_1$ is independent of $H$.  To see this explicitly, the
constant term lies in $[0,1)$, and an error coefficient at degree
$e$ contributes at most $A_1Z^{-e}(CZ)^e$ to the first bound and at
most $eA_1Z^{-e}(CZ)^{e-1}$ to a derivative.  Only boundedly many
monomials occur.

Choose a rational vector $V\in\mathfrak g$ with
$\xi(\exp V)=1$.  Such a vector exists because a nonzero horizontal
character has nonzero rational differential.  There is a fixed
positive integer $K$ with $\exp(KV)\in\Gamma$.  For completeness,
the Mal'cev coordinates of $\exp(tV)$ are rational polynomials in
$t$ with zero constant term, by the finite BCH formula.  Taking $K$
divisible by their finitely many denominators makes these coordinates
integral, which places $\exp(KV)$ in the lattice.  All choices depend
only on the fixed character and rational group data.

Put $H_\xi=\ker\xi$ and
\begin{equation}
 Q_z(x)=\exp(-a_z(x)V)P_z(x)\exp(-m_z(x)V).
 \label{eq:rough-kernel-polynomial}
\end{equation}
Then $\xi(Q_z(x))=0$ identically, for real $(z,x)$, and
\begin{equation}
 P_z(x)=\exp(a_z(x)V)Q_z(x)\exp(m_z(x)V).
 \label{eq:rough-kernel-factorization}
\end{equation}
The group $H_\xi$ is connected and simply connected: the exponential
map identifies it with the kernel of the linear differential of
$\xi$.  That kernel is rational, so $H_\xi$ is a rational subgroup
of dimension $\dim G-1$.  Formula~\eqref{eq:rough-kernel-polynomial}
and the finite BCH formula show that $\log Q_z(x)$ is jointly
polynomial in $(z,x)$ with a degree bound depending only on the old
degree and nilpotence step.  No bound on its coefficients is needed.

The polynomial now takes values in a smaller group.  To apply the
induction hypothesis, we must also express both sampling laws and
their test functions on a fixed compact quotient of $H_\xi$.
Partition each test box into a bounded number of rectangular
subboxes, each of side between $\lambda cZ/2$ and $2\lambda CZ$,
where the sufficiently small constant $\lambda>0$ is chosen below.
For example divide each side into the same sufficiently large fixed
number of intervals.  On a subbox choose a point $x_*$ and replace the
left factor $\exp(a_z(x)V)$ in
\eqref{eq:rough-kernel-factorization} by
$\ell_{z,*}=\exp(a_z(x_*)V)$.  The error in the test is at most
$C_2\lambda$: the derivative estimate
\eqref{eq:rough-slow-factor} bounds the change in $a_z$, and the action
map on the product of a fixed compact subset of $G$ and $G/\Gamma$
has bounded derivative.  The constant $C_2$ depends on the original
Lipschitz bound, but not on $H,z$ or the choice of subbox.  Choose
$\lambda$ so that twice this error is less than $\eta/8$.

Split each subbox further according to $x\bmod K$.  At a fixed $z$
and residue vector $u$, the integer polynomial $m_z(x)$ has a fixed
value $j\bmod K$, and hence
\[
 \exp(m_z(x)V)\Gamma=\sigma_j\Gamma,
 \qquad\sigma_j=\exp(jV),\quad0\leq j<K.
\]
The resulting test on the smaller group is
\begin{equation}
 h\bigl(H_\xi\cap\sigma_j\Gamma\sigma_j^{-1}\bigr)
       \longmapsto F_z(\ell_{z,*}h\sigma_j\Gamma).
 \label{eq:rough-kernel-test}
\end{equation}
The intersection is a lattice in $H_\xi$: conjugation by rational
$\sigma_j$ preserves the rational structure, and a connected rational
subgroup meets a rational lattice in a lattice.  This can also be
seen by taking a rational Mal'cev basis of the subgroup and clearing
the finite coordinate denominators.  The displayed map from its
compact quotient to $G/\Gamma$ is smooth.  Since there are only
finitely many $j$ and $\ell_{z,*}$ lies in a fixed compact set,
all the tests \eqref{eq:rough-kernel-test} have one uniform Lipschitz
bound.

We justify the sampling weights in this splitting.  For an interval
of length comparable to $Z$ and any residue modulo $q$, its count is
its length divided by $q$, with error at most two.  Multiplying this
formula over the fixed number of coordinates shows that the normalized
weights of the boundedly many subboxes under the unrestricted and
step-$t_z$ samples differ in total by $O(t_z/Z)$.  For $w\geq K$,
$\gcd(t_z,K)=1$.  Thus both samples assign weight
$K^{-d}+O(K^dt_z/Z)$ to each $K$-residue vector inside a subbox:
in the step-$t_z$ coordinate, reduction modulo $K$ is a bijection.
The total difference of the weights of all the pieces is consequently
$o(1)$, uniformly in $z$ and in the box endpoints.

After freezing, the original discrepancy is still at least
$7\eta/8$.  Write the two averages as convex combinations over these
pieces and replace one set of combination weights by the other.
The error is $o(1)$ times the uniform supremum bound on the tests.
It follows from the triangle inequality that at least one piece has
a conditional discrepancy at least $\eta/2$, for large $w$.
This lower bound does not depend on the number of pieces: it follows
from a convex combination, not from an unnormalized sum.

For each $z$ choose such a piece.  Its subbox and frozen observable may
vary arbitrarily with $z$, as allowed in the induction assertion.
Thin once more to make its residue vector $u\bmod K$ and its coset
index $j$ common.  In this subbox substitute
$x=u+Ky$, taking $u\in\{0,\ldots,K-1\}^d$.  The unrestricted
sample becomes all integer $y$ in a box of side comparable to $Z/K$,
within a fixed multiple of that scale.  The step-$t_z$ sample becomes
one residue vector modulo $t_z$, because $K$ is invertible modulo
$t_z$.  The new map $Q_z(u+Ky)$ is still jointly logarithmically
polynomial, and $Z/K$ still dominates every fixed power of $S$.
We have therefore obtained a bad progression family on
\[
 H_\xi/\bigl(H_\xi\cap\sigma_j\Gamma\sigma_j^{-1}\bigr)
\]
of strictly smaller dimension.

All the new degree, box, observable and discrepancy bounds are
independent of the progression length.  Characters, rational
directions, residue vectors and cosets range over fixed finite lists;
as explained at the start, one such quotient works for unbounded
lengths.  The sole length-dependent quantities were interpolation
denominators and asymptotic thresholds.  Those denominators enter only
the unrestricted coefficients of the kernel polynomial.  At every
retained integer index its spatial integer part has integer
coefficients, so the fixed residue splitting and the finite list of
quotient lattices are unchanged.  This contradicts the
induction hypothesis and proves that bad progression families do not
exist.  Failure of the Proposition produced precisely such a family,
so the Proposition follows, with the uniformity asserted in its
statement.
\end{proof}

\subsection{Residual steps and conditional sampling}

We record explicitly the form needed when the step being removed is
one factor of a larger product.  This separates the polynomial
requirement from the arbitrary choice of actual sampling box.

\begin{corollary}[Removing one factor while retaining the others]
\label{cor:rough-residual-step}
In Proposition~\ref{prop:rough-step}, let $q$ be another positive
integer, fixed while $t$ varies but permitted to change with $w$.
Let $a\in\{0,\ldots,q-1\}^d$ also be independent of $t$.
Suppose the spatial boxes have common scale $Z_0$, with
$Z_0/q$ dominating every fixed power of $S$.  Suppose
$\log P_t(a+qy)$ is polynomial jointly in $(t,y)$ with fixed degree
bound, and the normalized boxes lie in fixed multiples of
$[-Z_0/q,Z_0/q]^d$ with comparable sides.  Then, outside an exceptional
proportion tending to zero, averaging $P_t$ over the class
$a+q\Z^d$ agrees to any prescribed accuracy with averaging over any
compatible class modulo $qt$, with the same uniformities in boxes and
tests.  No coprimality between $q$ and $t$ is assumed.
\end{corollary}

\begin{proof}
Write the finer class as $b+qt\Z^d$ with $b\equiv a\pmod q$.
Under $x=a+qy$, its condition is exactly
\[
 y\equiv(b-a)/q\pmod t.
\]
This identity does not use invertibility of either factor.  Apply
Proposition~\ref{prop:rough-step} to the normalized polynomial and
the scale $Z_0/q$.  Its uniformity permits arbitrary compatible $b$
and arbitrary endpoints of the normalized boxes.
\end{proof}

In applying Corollary~\ref{cor:rough-residual-step}, the retained residue
class must be chosen independently of the factor being removed.
Section~\ref{sec:alignment} verifies this condition for each factor of
the tail product.

Finally, uniformity over additional conditioned parameters has a
precise probabilistic consequence.  Suppose $k$ factors range in fixed
dyadic intervals and smooth residue classes, and let $\nu$ be the
product of their normalized uniform laws.  Suppose a conditional law
$\mu$ satisfies $\mu\leq C_0\nu$, with fixed $C_0$.  If, for every
choice of the other $k-1$ factors, the exceptional proportion for
removing the remaining factor is at most $\epsilon_w\to0$, then
Fubini's Theorem gives $\nu(E_\ell)\leq\epsilon_w$ and hence
\[
 \mu\left(\bigcup_{\ell=1}^k E_\ell\right)
       \leq C_0k\epsilon_w=o(1).
\]
The uniform bound $\epsilon_w$ follows from
Proposition~\ref{prop:rough-step}: otherwise choices of the additional
parameters with exceptional proportion bounded below would themselves
form a violating sequence of permitted data.  This gives the conditional
form of the progression comparison used in Section~\ref{sec:alignment}.
